\section*{Chapitre \ref{ch:m1:exchanges-brokers-venues} --- Bourses, courtiers et plateformes de négociation}

\begin{solution}{exo:m1:exchanges-brokers-venues:1}
Sur un total de 5\,411 millions de dollars~: énergie 40.3~\%~; actions au comptant et options sur actions 8.6~\%~;
données et connectivité 19.1~\%~; admissions à la cote 9.1~\%.
\end{solution}

\begin{solution}{exo:m1:exchanges-brokers-venues:2}
$2 \times 0.70/280\,000 = 5\times10^{-6}$~: \qty{0.05}{\bp}. Une transaction sur action
portant sur la même valeur coûterait nettement plus cher, rien que par la fourchette.
\end{solution}

\begin{solution}{exo:m1:exchanges-brokers-venues:3}
(a) Algorithme du \omterm{def:m1:exchanges-brokers-venues:broker}{courtier} ou desk~; contrôles du \omterm{def:m1:exchanges-brokers-venues:broker}{courtier}. (b) \omterm{def:m1:exchanges-brokers-venues:dma}{Accès direct au
marché} (DMA)~; contrôles du \omterm{def:m1:exchanges-brokers-venues:broker}{courtier}, dans ses systèmes. (c) \omterm{def:m1:exchanges-brokers-venues:dma}{Accès sponsorisé}~;
toujours les contrôles du \omterm{def:m1:exchanges-brokers-venues:broker}{courtier}, sous son contrôle exclusif, bien qu'ils tournent
aux côtés des machines du client. (d) Membre de la \omterm{def:m1:exchanges-brokers-venues:exchange}{bourse}~; les contrôles propres de la société, dont
elle répond auprès de son régulateur et de son membre compensateur.
\end{solution}

\begin{solution}{exo:m1:exchanges-brokers-venues:4}
180\,000~\$~; 65\,000~\$~; 49\,000~\$~; 150\,000~\$. L'\omterm{def:m1:exchanges-brokers-venues:dma}{accès sponsorisé} permet
d'économiser 16\,000~\$ par mois par rapport au DMA. Le statut de membre l'emporte au-delà de 312.5 millions
d'actions par mois, soit cinq fois le volume de cette société.
\end{solution}

\begin{solution}{exo:m1:exchanges-brokers-venues:5}
Du DMA vers le sponsorisé~: $20\,000/(0.0010-0.0006) = 50$ millions d'actions par mois
(contre 33). Du sponsorisé vers le statut de membre~: $125\,000/0.0006 = 208$ millions (contre
313). La plage où l'\omterm{def:m1:exchanges-brokers-venues:dma}{accès sponsorisé} est le plus avantageux se réduit par ses deux extrémités~: le
fournisseur perd ses plus petits clients au profit du DMA et pousse les plus gros à devenir
membres --- ceux qui le rémunéraient le plus.
\end{solution}

\begin{solution}{exo:m1:exchanges-brokers-venues:6}
$x^\star = 40/(504+50) = 7.2\%$~; à 10~\%, $0.10 \times 554 - 40 = 15.4$
millions de dollars, contre 40.6 millions de dollars avant. Une baisse d'un centième de cent
fait disparaître les trois cinquièmes du profit.
\end{solution}

\begin{solution}{exo:m1:exchanges-brokers-venues:7}
Seuls l'\omterm{def:m1:exchanges-brokers-venues:dma}{accès sponsorisé} (50 microsecondes) et le statut de membre (40) respectent 1
milliseconde. À 5 millions d'actions par mois, l'\omterm{def:m1:exchanges-brokers-venues:dma}{accès sponsorisé} coûte 27\,000~\$
contre 150\,000~\$. Sans contrainte, le choix serait le DMA à 10\,000~\$~: la
contrainte de latence coûte 17\,000~\$ par mois, soit 0.34 cent par action --- davantage que
ce que facturait l'algorithme du \omterm{def:m1:exchanges-brokers-venues:broker}{courtier} pour tout.
\end{solution}

\begin{solution}{exo:m1:exchanges-brokers-venues:8}
Le projet compare les frais et laisse de côté le \emph{prix} obtenu et la
probabilité de l'obtenir. Le coût par action d'un preneur de liquidité est égal aux frais augmentés de la
demi-fourchette disponible sur la plateforme, et un ordre qui ne peut pas y être exécuté
coûte le délai et le mouvement du prix dans l'intervalle. La comparaison consiste à
comparer $\text{frais}_A + s_A/2$ contre $\text{frais}_B + s_B/2$, pondéré par
la probabilité d'exécution. La plateforme établie dispose des ordres, de sorte que sa fourchette vaut un cent
là où le carnet d'ordres du nouvel entrant est vide ou large de cinq cents~: une économie de 0.25 cent
sur les frais, contre 2 cents de fourchette en plus. Sans liquidité, la plateforme bon marché est
la plus chère, et un faible tarif pour les preneurs de liquidité ne fait rien pour attirer les ordres au repos
qui changeraient cette situation.
\end{solution}

\begin{solution}{pb:m1:exchanges-brokers-venues:1}
\textbf{1.} $0.28 - 0.25 = 0.03$ cent.
\textbf{2.} $0.0003/45 = \qty{0.067}{\bp}$.
\textbf{3.} $8\times10^9 \times 252 \times 0.0003 = 604.8$ millions de dollars~; 6.05
millions de dollars par point.
\textbf{4.} $6.05 + 0.60 = 6.65$ millions de dollars par point.
\textbf{5.} $38/664.8 = 5.72\%$.
\textbf{6.} 457 millions d'actions, 20.6 milliards de dollars par jour.
\textbf{7.} $12\% \times \tfrac13 \times \tfrac12 = 2\%$.
\textbf{8.} Non~: $2 \times 6.65 - 38 = -24.7$ millions de dollars par an.
\textbf{9.} $\nu = -0.04$ cent~: la plateforme rémunère l'appariement. À 4~\%~:
$0.04 \times 8\times10^9 \times 252 \times (-0.0004) = -32.3$ millions de dollars.
\textbf{10.} $-32.3 + 2.4 - 38 = -67.9$ millions de dollars.
\textbf{11.} $6 \times 6.65 - 38 = 1.9$ millions de dollars~: trente-six années de ce type pour
rembourser une année de subvention. À 6~\%, la plateforme survit, mais elle ne vaut pas grand-chose.
\textbf{12.} $38/(403.2 + 60) = 8.2\%$.
\textbf{13.} En tant que propriétaires, ils se versent à eux-mêmes une partie des frais qu'ils paieraient
de toute façon, et --- la vraie raison --- une alternative crédible discipline les
frais de l'acteur en place sur les 94~\% restants de leur volume. Une plateforme qui perd
25~millions de dollars, mais réduit de 0.02 cent les frais de l'acteur en place sur le flux du consortium
peut s'autofinancer.
\textbf{14.} Elle renonce au droit d'admission à la cote, à sa part des revenus de données consolidées
et au statut réglementaire qui lui permet d'édicter des règles contraignantes~; elle
économise le coût et les délais de l'agrément de \omterm{def:m1:exchanges-brokers-venues:exchange}{bourse} et de l'autorégulation.
\textbf{15.} Elle attire les ordres au repos des institutions plus lentes et des \omterm{def:m1:what-a-trading-firm-does:market-maker}{teneurs de
marché} qui redoutent de se faire prendre à contre-pied, le délai leur permettant de recoter
en premier~; elle repousse les preneurs les plus rapides, dont l'avantage tient précisément à cette course.
\textbf{16.} Le bénéfice que chaque trader tire d'une plateforme croît avec les ordres déjà
présents, de sorte que la part de marché s'auto-alimente~: quelques plateformes, chacune dominante dans une niche
(un modèle de frais, un type d'ordre, un groupe de clients), survivent~; les autres restent en deçà de
$x^\star$ et ferment ou sont rachetées.
\textbf{17.} Si les revenus de données sont répartis en partie selon le temps passé à coter au
meilleur cours, les ordres placés au meilleur cours mais annulés avant
d'avoir pu être exécutés rapportent des revenus à la plateforme sans fournir de liquidité~; la
plateforme peut les rémunérer par des paliers de rabais.
\textbf{18.} Que chaque sponsor soit membre en règle d'un dispositif de compensation,
qu'il dispose de contrôles pré-négociation sur le flux sponsorisé, sous son contrôle exclusif,
et qu'il puisse couper l'accès du client~; la plateforme applique aussi ses propres
limites au niveau de l'ordre et un coupe-circuit par identifiant.
\textbf{19.} \textbf{20.6 milliards de dollars par jour.}
\textbf{20.} Réglementaire et commercial~: la technologie d'appariement est une
marchandise banalisée, alors que la licence, les droits sur les données et surtout le flux d'ordres
engagé par ses propriétaires décident de sa survie.
\end{solution}

\begin{solution}{iq:m1:exchanges-brokers-venues:1}
Frais d'admission à la cote, frais de transaction nets des rabais, frais de compensation et revenus sur
le collatéral lorsqu'elle possède la chambre de compensation, licences de données de marché, et accès
(ports, colocation, connectivité). Pour une société à haute fréquence, la ligne transactions
est souvent proche de zéro, voire négative après rabais~; ce pour quoi elle paie est
les données et l'accès, la partie fixe et croissante.
\iqlookfor{les cinq postes, et la conscience que la facture de la société à haute fréquence porte sur
les données et la connectivité.}
\end{solution}

\begin{solution}{iq:m1:exchanges-brokers-venues:2}
Les ordres DMA transitent par les systèmes du \omterm{def:m1:exchanges-brokers-venues:broker}{courtier}~; les ordres en \omterm{def:m1:exchanges-brokers-venues:dma}{accès sponsorisé} vont
directement du client à la \omterm{def:m1:exchanges-brokers-venues:exchange}{bourse}, sous l'identifiant du \omterm{def:m1:exchanges-brokers-venues:broker}{courtier}. Dans
les deux cas, le \omterm{def:m1:exchanges-brokers-venues:broker}{courtier} répond devant la \omterm{def:m1:exchanges-brokers-venues:exchange}{bourse} et le régulateur~: depuis la
règle sur l'accès au marché, il doit exercer des contrôles préalables sur ce
flux, sous son contrôle direct et exclusif, si bien que les ordres erronés sont
une défaillance de son propre contrôle, autant que le bug du client --- et son argent si le client ne
peut pas payer.
\iqlookfor{la responsabilité suit l'identifiant de membre.}
\end{solution}

\begin{solution}{iq:m1:exchanges-brokers-venues:3}
Par ordre~: quantité et notionnel maximaux~; bande de prix par rapport à un prix de
référence~; instrument autorisé et négociable~; indicateurs de vente à découvert et réglementaires.
Au niveau agrégé~: limites de position et de notionnel par instrument et globales, exposition
aux ordres en attente, cadences d'ordres et de messages, limite de perte, protection contre les doublons et
de l'autoappariement, plus un coupe-circuit. Le budget se mesure en microsecondes ou
moins, si bien que les contrôles sont des calculs en mémoire, peu ramifiés, sur un état
tenu à jour de façon incrémentale --- pas de verrou, pas d'allocation, aucun appel à une base de données ou au réseau
dans le chemin critique --- et leurs limites sont chargées au démarrage et ne sont modifiées que par
un canal auxiliaire audité.
\iqlookfor{contrôles agrégés aussi bien que contrôles par ordre, et conséquences du budget pour
la conception.}
\end{solution}

\begin{solution}{iq:m1:exchanges-brokers-venues:4}
Une action est fongible d'une plateforme à l'autre et la loi oblige les plateformes à se disputer
ce flux, si bien que les frais d'appariement ont glissé vers le coût. Un contrat à terme est le produit
propre de la \omterm{def:m1:exchanges-brokers-venues:exchange}{bourse} et se compense dans sa chambre de compensation~: les positions ouvertes ne peuvent
pas migrer, donc la liquidité non plus. La \omterm{def:m1:exchanges-brokers-venues:exchange}{bourse} de contrats à terme est un monopole sur chaque
contrat qui réussit, et elle tarife comme tel.
\iqlookfor{fongibilité et compensation comme mécanisme, et non «~les contrats à terme sont
plus gros~».}
\end{solution}

\begin{solution}{iq:m1:exchanges-brokers-venues:5}
Comparer le profit espéré par unité de temps, et non le rabais par exécution~: taux d'exécution
$\times$ (demi-fourchette $+$ rabais $-$ \omterm{def:m1:what-a-trading-firm-does:adverse-selection}{sélection adverse}). Un cinquième des
exécutions exige cinq fois la \omterm{prop:m1:what-a-trading-firm-does:capture}{captation nette} pour rivaliser, et un rabais de 0.12 cent
plus élevé ne suffit pas --- sauf si les exécutions sur la nouvelle plateforme sont beaucoup
moins toxiques, ce qui est une question mesurable (les markouts par plateforme). Coter
des deux côtés si le capital et les limites de risque le permettent~; le rabais peut aussi compter
dans un palier de volume.
\iqlookfor{markouts par plateforme, et raisonner en taux plutôt qu'en économie
par exécution.}
\end{solution}

\begin{solution}{iq:m1:exchanges-brokers-venues:6}
Équitable~: les produits sont proposés à tout le monde à un prix publié~; ceux qui les achètent
sont ceux qui maintiennent des cotations serrées et cohérentes d'une plateforme à l'autre, et
les investisseurs obtiennent de meilleurs prix en conséquence~; la vitesse est un coût de cette activité
comme une autre. Inéquitable~: la \omterm{def:m1:exchanges-brokers-venues:exchange}{bourse} vend à certains un avantage sur ses autres clients,
avantage qui n'existe que parce qu'elle publie aussi une version plus
lente à leur intention~; les frais sont fixés par un monopole sur ses propres données~; et la
course qui en résulte consomme des ressources à grappiller des microsecondes, sans profit pour
les investisseurs finaux, alors que c'est au flux lent que se mesurent les obligations de
meilleure exécution.
\iqlookfor{les deux thèses défendues correctement~; le point de vue propre du candidat compte moins
que de voir la structure.}
\end{solution}
